\section{验证}

\begin{tcblisting}{consolebox}
# 串口输出：
[init] gdt ok
\end{tcblisting}

如果 GDT 有误（如 access byte 写错、base/limit 计算错误），CPU 会在
far jump 时 Triple Fault 并重启。

\subsection{GDB 逐指令验证}

\begin{tcblisting}{consolebox}
# 终端 A
«\tps{\$}» «\tin{make DEBUG=1 debug}»

# 终端 B
«\tps{\$}» «\tin{make gdb}»
«\tps{(gdb)}» «\tin{b gdt\_flush}»
«\tps{(gdb)}» «\tin{c}»
Breakpoint 1, gdt_flush () at gdt_flush.asm:21
«\tps{(gdb)}» «\tin{info reg cs ds ss}»
cs  0x8   8           # GRUB 设置的 CS（通常也是 0x08）
ds  0x10  16          # GRUB 设置的 DS
ss  0x10  16

«\tps{(gdb)}» «\tin{si}»              # lgdt [eax]
«\tps{(gdb)}» «\tin{si}»              # jmp 0x08:.flush_cs
«\tps{(gdb)}» «\tin{info reg cs}»
cs  0x8   8           # CS 重新加载，hidden cache 已更新

«\tps{(gdb)}» «\tin{si}»              # mov ax, 0x10
«\tps{(gdb)}» «\tin{si}»              # mov ds, ax
«\tps{(gdb)}» «\tin{info reg ds}»
ds  0x10  16          # DS 重新加载

«\tps{(gdb)}» «\tin{c}»               # 继续，应该看到串口打印 [init] gdt ok
\end{tcblisting}

% ============================================================================

